\subsection{Equal amounts still have separate source occurrences}

Suppose an issuer repeats one tagged amount in a statement and a note. A query
may need a single amount, but an auditor still needs both locations and the
evidence that they agree. These are different requirements. The proposed
exchange preserves each occurrence before constructing a view that groups
facts for a particular question. The inspected Arelle methods make the reason
for that separation precise: displaying a value, comparing two values,
grouping facts with matching aspects and removing repeated facts perform
different operations \citep{arellefactidentity2026}.

A displayed value is particularly unsuitable as the retained amount. The
\code{effectiveValue} property produces a view: it rounds and formats numbers,
represents nil with a string, and can return an error description. Even a
convenient numeric property needs scrutiny. The inspected \code{vEqValue}
converts numeric text to a binary floating-point number. The distinct integers
$2^{53}=9{,}007{,}199{,}254{,}740{,}992$ and $2^{53}+1$ then have the same
floating-point representation. This constructed example does not describe an
issuer error; it shows why that conversion cannot establish exact identity.
The adapter therefore keeps the source literal, the typed value and the
displayed reading as distinct fields, with the transformation between them
recorded \citep{arellefactidentity2026}.

The same caution applies to a method whose name contains \emph{equal}. In the
inspected \code{isVEqualTo}, the ordinary numeric comparison checks context
and unit and follows the selected rounding mode, but the method does not itself
compare concept names. Its nil branch returns whether the other fact is nil
before comparing contexts. A caller needing facts about the same accounting
quantity must establish that condition separately. Conversely, the item branch
of \code{isDuplicateOf} checks concept, parent, context and numeric unit
without comparing the item values. It can identify candidates for investigating
a disagreement; its result does not resolve that disagreement. The tuple
branch has additional child comparisons and must not be described as the same
test \citep{arellefactidentity2026}.

The dedicated duplicate-fact module supplies more explicit grouping. It checks
concept, context and unit after forming hash buckets, and includes language
for the relevant text-fact types. Its value tests nevertheless remain specific
to the operation. Consider two hypothetical valid numeric occurrences with the
same concept, context, unit and declared decimals of zero, whose value strings
are \code{1} and \code{1.0}. In the inspected code, the ordinary value-equality
key retains those strings, so the all-complete classification distinguishes
them. The complete-subset reduction instead keys on declared decimals and the
typed value; decimal one equals decimal one, so it retains the first occurrence.
Thus even within this module, the classification and reduction do not share
one universal equality key. The source strings and membership of a reduced
view must remain recoverable \citep{arellefactidentity2026}.

Reported accuracy creates a second distinction. For a finite declared decimal
accuracy $d$, the inspected non-truncating interval helper assigns a numeric
value $v$ the inclusive bounds $v-\tfrac12 10^{-d}$ and
$v+\tfrac12 10^{-d}$. These are the helper's comparison bounds, not a confidence
interval for the true economic amount. Values zero and one, both declared to
zero decimals, have bounds $[-0.5,0.5]$ and $[0.5,1.5]$. They touch. The
interval option of \code{isVEqualTo} therefore admits overlap under its other
required conditions. The duplicate-set consistency check rejects the pair
because unequal typed values with the same declared decimal accuracy fail its
additional condition. Substituting one test for the other would change the
question being answered \citep{arellefactidentity2026}.

Pairwise agreement also differs from agreement of the whole set. Take three
constructed observations of the same quantity: 100 at zero decimals, and 99.6
and 100.4 at one decimal. The first has bounds $[99.5,100.5]$; the other two
have bounds $[99.55,99.65]$ and $[100.35,100.45]$. Each precise observation
overlaps the coarse one, but the precise observations disagree and the three
intervals have no common point. The inspected consistent-pairs reduction can
remove the coarse occurrence while leaving both precise occurrences. The
consistent-set path instead reports the inconsistent set and falls back to
complete-subset reduction. Neither operation establishes which precise
observation the accounting analysis should accept \citep{arellefactidentity2026}.

These paths require their own admission checks. The initial grouping considers
validation state and the availability of a concept, type and context; a
singleton bucket is not itself a validity verdict. The numeric consistency
path also expects an inferred decimal accuracy from its helper, which returns
no value when the fact's decimals attribute is absent. An adapter must not
silently reinterpret an unsupported or unvalidated case as an ordinary
consistent numeric fact. It retains the declared decimals and precision,
validation state, nil state, concept, full context and unit, together with the
particular comparison operation and its options \citep{arellefactidentity2026}.

Finally, the inspected deduplicated-instance export removes fact nodes from
the model's underlying tree before saving. The controller performs it at the
end and warns that subsequent use of the modified model can be unreliable.
The proposed exchange therefore copies every original occurrence and its
diagnostics before constructing any optional reduced view. A view records
its input members and selection rule; it does not replace the source record.
This preserves disagreement for review while allowing a query to avoid
counting an accepted repeated amount twice. These are requirements derived
from pinned source inspection and independently checked constructed arithmetic;
an executed adapter and actual-filing validation remain to be established
\citep{arellefactidentity2026}.
